\subsection{二中心fiberによる固定楕円パラメータの全排除}
\label{subsec:cubic-third-root-unit-boundary-L-open-obstruction}
\begin{proposition}[全変形の相対Frobenius双対]
$E=2$ の全 $u,v$ に対し
\[
 \operatorname{Hom}_{k[x,\Psi,z_9]}(A_{u,v},k[x,\Psi,z_9])
 \simeq (A_{u,v})_\nu(12),\qquad
 \nu(x,y,b,c)=(x,y,\omega^2b,\omega c).
\]
\end{proposition}
\begin{proof}
パラメータ独立なArtinian商の語基底81個を用い、固定例と同じtop functionalを定める。
その完全pairingはArtinian商の全13次数で直接確認する。
相対版の完全性は次数付きNakayamaによる。
残る双対の自己同型については、$A/(\Psi,z_9)$ を
$K[u,v]$ 上で次数13まで構成する。
$x$ 倍と自由基底変換の二つの $81\times81$ 行列は、
逆行列の共通分母がパラメータ定数であることを確認して反転する。
全81基底元と四生成元の324等式を多項式として再代入し零を確認する。
再現器は \texttt{verify\_cubic\_third\_root\_unit\_boundary\_L\_polynomial\_Frobenius.py}。
\end{proof}

\begin{lemma}[trace非退化開集合上の中心]
\label{lem:cubic-trace-open-center}
$k$ は標数0代数閉体、$P$ は既約な被約有限型 $k$ 代数、$V$ は $P$ 上有限自由な代数とする。
正則trace形式が非退化な開集合が非空で、その一点でfiberの中心が一次元なら、
その開集合の全fiberは中心単純である。
\end{lemma}
\begin{proof}
trace非退化な有限次元代数ではJacobson radicalは零であり、各fiberは半単純となる。
trace双対基底 $e_i,e^i$ を用いて
\[
 E(a)=\sum_i e_i a e^i
\]
を定める。各半単純fiberのmatrix blockでは、正則traceの規格化から
$E$ は中心へのprojectorとなる。
従って $E^2=E$、像が全中心となる等式は、被約な開集合全体でも成り立つ。
このprojectorのrankは局所一定であり、開集合は既約なので一定。
既知の一点でrank1であるから、全fiberの中心は一次元である。
\end{proof}

\begin{theorem}[$E=2$ の全二変形で中心単純次数9]
\label{thm:cubic-third-root-unit-boundary-L-E2-CSA}
$E=2$、$L=0$ の全六埋込みに対し、$v\ne0$ の全 $u$ で
$A_{u,v}$ の全中心は $k[x,\Psi,z_9]$、一般点は次数9の中心単純代数である。
\end{theorem}
\begin{proof}
$x$ のscalar変換により $v=1,u=a$ に正規化する。
$P=K[a]$ 上でもArtinian基底の持上げは
$P[x,\Psi,z_9]$ 上の自由基底になる。
実際、三中心元による商はパラメータ独立なArtinian商の $P$ への拡大であり、
次数付き帰納法で基底持上げからの写像は全射となる。
その核は自由加群の部分加群で一般点で零なので零である。

$u=1,v=0$ における新しい $z_9$ を、既知の旧中心基底で書く：
\[
 z_9=a_0x^9+a_1x^6\Psi+a_2x^3\Psi^2+a_3\Psi^3+c z_9^{\rm old},\qquad c\ne0.
\]
係数は公開証明書の中心基底変換で一意に定まり、元の全語ベクトルへ再代入して確認する。
$b=a_0+c$ とおく。
次の二つの $P$ 上自由階数81の中心fiberを使う：
\[
 V_b=A_{a,1}/(x-1,\Psi,z_9-b),\qquad
 V_1=A_{a,1}/(x-1,\Psi,z_9-1).
\]
それぞれの正則trace Gram行列式を $f_b(a),f_1(a)\in K[a]$ とする。
各々、$a=0$ で中心一次元かつtrace非退化となることを正の有限体witnessで確認する。
補題\ref{lem:cubic-trace-open-center}により、各 $D(f_\bullet)$ の全fiberは中心単純である。

従って $f_b,f_1$ に共通零点がないことを示せばよい。
その証明では、先に標数0の全乗法行列を構成し、その後で還元する。
$A/(\Psi,z_9-\beta x^9)$ の次数12,13の語基底と $x$ 倍の逆行列を用いると、
$v=1,u=a$ の全生成元行列は $K[a]$ の多項式となる。
全relation pivotと $x$ の逆行列分母がパラメータ定数であることを確認する。
基底語の $y,b,c$ の長さを $d_i$ とすれば、全生成元行列について
\[
 2\deg_a (R_v)_{ji}\le d_i+1-d_j
\]
を標数0の実際の全係数で確認できる。
従って $\deg_a\operatorname{Tr}(R_{e_i e_j})\le(d_i+d_j)/2$、
\[
 \deg f_b,\deg f_1\le\sum_i d_i=486.
\]

還元 $p=1021,t=169$ を使う。
六次最小多項式、原始三乗根条件、全係数の有理分母が $p$ で可逆なことを確認する。
各fiberのtrace行列式を487点で評価し、次数上界486を用いて一意に補間する。
さらに8点の再評価を確認する。
得られる二式は
\[
 \overline f_b=340h_b^{54},\qquad \overline f_1=340h_1^{54},
\]
ただし $\mathbb F_{1021}[a]$ において
\begin{align*}
 h_b={}&(a^2+228a+306)(a^2+471a+147)(a^2+573a+1012)\\
 &\hspace{12mm}\cdot(a^3+317a^2+112a+981),\\
 h_1={}&(a+95)(a+110)(a+579)\\
 &\hspace{12mm}\cdot(a^6+805a^5+322a^4+741a^3+284a^2+740a+787).
\end{align*}
全係数を掛け戻して補間式との一致を確認し、Euclid算法で
$h_b,h_1$ のB\'ezout恒等式1を確認する。
二つの486次trace多項式のresultantは $965\bmod1021$ で非零である。

両還元の次数が上界486に達するので、標数0でも次数は正確に486であり、
次数落ちした還元のgcdだけを見ているわけではない。
同じ固定次数のSylvester resultantが非零に還元されるから、
$f_b,f_1$ の標数0でのresultantも非零である。
従って代数閉体上でも共通零点はない。
全六埋込みに対して、各 $a$ で少なくとも一方の中心fiberが中心単純となる。

その正のtrace行列式と中心minorを中心一般点へ移せば、
$A_{a,1}$ の一般点は次数9の中心単純代数である。
中心部分環の整閉性から全中心は $k[x,\Psi,z_9]$ に一致する。
再現器は \texttt{verify\_cubic\_third\_root\_unit\_boundary\_L\_two\_fibers.py}。
NumPyの整数行列演算では全積の直後に還元し、
内積上限 $81\cdot1020^2=84272400<2^{63}$ を確認する。
これは有限体点の標本からの外挿ではなく、標数0次数上界と
非零resultantを組み合わせた証明である。
\end{proof}

\begin{theorem}[$E=2$ の正常一次境界の全H3排除]
\label{thm:cubic-third-root-unit-boundary-L-E2-no-H3}
$E=2$、$L=0$ の全六埋込みで、$w=1$ の全 $Z,u,v\in k$ はH3を与えない。
\end{theorem}
\begin{proof}
命題\ref{prop:cubic-third-root-unit-boundary-L-incidence}のshearにより
$Z=0$ へ移せる。この操作は $d,E$ を変えず、$u,v$ だけを変える。
従って $Z=0$ の全 $u,v$ を排除すればよい。
$v\ne0$ では前定理から全中心が多項式環、一般中心単純次数が9である。
相対双対は位数3の $\nu$ によるものなので、
有限大域次元なら局所order認識と周期障害によりBrauer周期は3を割る。
一方、整域なら曲面period--index定理により周期9となる。
従ってH3の二つの必要条件は両立しない。
$v=0$ は既にscalar同型と頂点障害で全 $u$ を排除した。
\end{proof}

\begin{remark}[独立な第二点と適用範囲]
$u=0,v=1$ の中心fiber $\Psi=0,z_9=b$ では、
$p=109,t=15,\omega=63$ の還元で $b=83$、中心一次元、
正則trace Gram行列式 $64\bmod109$ を独立に確認している。
再現器 \texttt{verify\_cubic\_third\_root\_unit\_boundary\_L\_second\_CSA.py}
は元の六二次関係、三次切断、二中心関係の全行列への再代入も確認する。
本節は固定楕円パラメータ $E=2$ の全変形を扱う。
全 $L$ 曲線の排除、各商の整域性や周期の単独決定、全H3分類は未完成である。
\end{remark}
